\documentclass[addpoints]{exam}
% \documentclass[addpoints,12pt]{exam}

\usepackage[slovene]{babel}
\usepackage{amsfonts}
\usepackage{amssymb}
\usepackage{amsmath}
\usepackage{float}
\usepackage{graphicx}
\usepackage{enumitem}
\usepackage{color}
\usepackage{eurosym}

\usepackage{pgfpages}
% \usepackage{tikz}
\usepackage{wrapfig}
\usepackage{pgfkeys}
\usepackage{pgfplots}
\usepackage{xcolor}
\usepackage{tkz-euclide}
\usepackage{xfp}




%Komentar naslednje vrstice glede na verzijo testa -- rešitve ali prostor za odgovore:
% \printanswers

% \linespread{2}


\pointsinrightmargin
\marginbonuspointname{}


\renewcommand{\solutiontitle}{\noindent\textbf{Rešitve in točkovanje:}\par\noindent}

\colorgrids
\definecolor{GridColor}{gray}{0.7}
\setlength{\gridsize}{7mm}

\extrawidth{5mm}
\setlength{\rightpointsmargin}{1.5cm}

\begin{document}

\runningfooter{}{}{}

\begin{center}
    \fbox{\fbox{\parbox{15cm}{\centering
    Drugo pisno ocenjevanje znanja -- ponovno \\ 2. letnik -- test E \\ 12. januar 2026 \\ (Kotne funkcije; Vektorji)}}}
\end{center}

\vspace{0.1in}
\makebox[\textwidth]{Ime in priimek, oddelek:\enspace\hrulefill}


\begin{center}
    \chqword{Naloga}
    \chpword{Možne točke}
    \chbpword{Dodatne točke}
    \chsword{Dosežene točke}
    \chtword{Skupaj}  
    % \combinedpointtable[h][questions]
    \combinedgradetable[h][questions]

\end{center}

\vspace{0.1in}


\begin{questions}

    % 1. vprašanje

    \question[4]
    Poenostavite izraz. 
    $$ \left(\sin^{-1}x\dfrac{1}{\cot x}\right)^{-2}+\sin x\tan x\cos x $$

    \begin{solution}[6.5cm]
       
    \end{solution}


    % 2. vprašanje

    \question[4]
    Izračunajte natančno vrednost izraza. 
    $$ \cos^2 55^\circ + \sin^2 305^\circ - \dfrac{\frac{\tan 213^\circ}{\sin 33^\circ}}{\sin 57^\circ} $$

    \begin{solution}[6cm]
       
    \end{solution}


\newpage
    % 3. vprašanje

    \question
    V kocki $ABCDEFGH$ točka $M$ tretjini rob $HG$, točka $N$ razdeli rob $AE$ v razmerju $|AE|:|NE|=5:1$.
    Z baznimi vektorji $\overrightarrow{a}=\overrightarrow{AB}, \overrightarrow{b}=\overrightarrow{AD}$ in $\overrightarrow{c}=\overrightarrow{AE}$ izrazite vektorje: 
    \begin{parts}
        \part[3]
        $\overrightarrow{MN}$, 

    \begin{solution}[5cm]

    \end{solution}
        
        \part[2] 
        $\overrightarrow{NB}-\overrightarrow{HB}+\overrightarrow{HA}-\overrightarrow{EA}$.
    \end{parts}

    \begin{solution}[5cm]

    \end{solution}

    

    
    % 4. vprašanje

    \question[4]
    V trikotniku $\triangle ABC$ z oglišči $A(4,5,-6)$, $B(1,7,0)$ in $C(4,0,-3)$ izračunajte oddaljenost težišča trikotnika $\triangle ABC$ od oglišča $B$.
    % \begin{parts}

    %     \part[2]
    %     dolžino težiščnice na stranico $AC=b$, 

    % \begin{solution}[6cm]

    % \end{solution}

    %     \part[3] 
    %     koordinate težišča trikotnika $\triangle ABC$.

    % \begin{solution}[2cm]

    % \end{solution}        
    
    % \part[5]
    %     skalarni produkt $\overrightarrow{AB}\cdot\overrightarrow{AC}$ in kot $\angle BAC=\alpha$, 

    %     \begin{solution}[5cm]

    % \end{solution}




    % \end{parts}







    \newpage
    % 5. vprašanje

    \question
    Daljica $AB$ je določena s krajiščema $A(7,5,-2)$ in $B(11,1,5)$.
    
    \begin{parts}
        \part[2]
        Izračunajte skalarni produkt krajevnih vektorjev točk $A$ in $B$.
       
           \begin{solution}[3cm]

    \end{solution}

        \part[4]    
        Določite enotski vektor v smeri krajevnega vektorja točke $S$, ki je razpolovišče daljice $AB$.

    \begin{solution}[6cm]

    \end{solution}

        \part[3]
        Določite $x$, da bo vektor $\overrightarrow{v}=(x-8,x,-7)$ vzporeden vektorju $\overrightarrow{AB}$.

            \begin{solution}[6cm]

    \end{solution}

        \part[3]
        Določite $y$, da bo vektor $\overrightarrow{u}=(y,6y+1,10)$ pravokoten na krajevni vektor točke $B$.
    \end{parts}

\begin{solution}[4cm]

\end{solution}


\newpage

    % 6. vprašanje



    \question[4]
    Dolžina vektorja $\overrightarrow{y}$ je $8$, dolžina vektorja $\overrightarrow{x}$ je $6$, 
    kot med njima pa meri $60^\circ$. 
    Izračunajte dolžino vektorja $\overrightarrow{z}=\frac{1}{4}\overrightarrow{y}+\frac{4}{3}\overrightarrow{x}$.
        


    \begin{solution}[10cm]

    \end{solution}





    \bonusquestion[3]
    \textit{Dodatna naloga:} Pokažite, da vektorja $\overrightarrow{a}=\left(\frac{3}{5},\frac{4}{5}\right)$ in $\overrightarrow{b}=\left(-\frac{4}{5},\frac{3}{5}\right)$ določata ortonormirano bazo ravnine.

    \begin{solution}[1cm]

    \end{solution}

\end{questions}



\end{document}